
\chapter{Summary of papers}

This chapter summarizes the three publications that are the foundation of this thesis. The unifying subject of this thesis is the investigation of transcription factor-DNA interactions. I explored this subject through two main aims:

    (1) discover new TF binding profiles and update the JASPAR database;

    (2) identify and characterize co-binding DNA patterns.

\section{Paper I. The release of the 2022 JASPAR database}

Castro-Mondragon JA\footnotemark[\value{footnote}], Riudavets-Puig R\footnotemark[\value{footnote}], \textbf{\underline{Rauluseviciute I}}\footnotemark[\value{footnote}], Berhanu Lemma R, Turchi L, Blanc-Mathieu R, Lucas J, Boddie P, Khan A, Manosalva Pérez N, Fornes O, Leung TY, Aguirre A, Hammal F, Schmelter D, Baranasic D, Ballester B, Sandelin A, Lenhard B, Vandepoele K, Wasserman WW, Parcy F, Mathelier A. \textit{JASPAR 2022: the 9th release of the open-access database of transcription factor binding profiles.} Nucleic acids research. 2022 Jan 7;50(D1):D165-73; doi: 10.1093/nar/gkab1113. \footnotetext{These authors should be regarded as joint first authors.}
\\

\textbf{Paper I} describes the 2022 release of the JASPAR database. JASPAR is a widely used open-access repository of transcription factor binding motifs created and maintained since 2004. The database is updated regularly with new motifs, features, and tools. One of the biggest strengths of JASPAR is that a team of experts manually curates all motifs before they enter the database. We manually curate two types of motifs: (1) \textit{de novo} discovered motifs from peak data and (2) external motifs from individual papers or databases. We mine peak data from databases (for example, ReMap \cite{hammal2022remap}, and GTRD \cite{kolmykov2021gtrd}) or from publications. Currently, JASPAR has two motif collections: CORE and UNVALIDATED. CORE collection contains profiles with orthogonal support. The UNVALIDATED collection has profiles of the same sound quality but without validation by independent experiments.

In 2022, CORE and UNVALIDATED collections were enhanced or updated by 501 and 182 profiles. The CORE collection increased by 19\% compared to the previous release. In addition to expanding the motif collections, we also started focusing on revising existing motifs. We revised UNVALIDATED collection profiles to check for possible new validations that would allow us to move the profiles into the CORE collection. In 2022, we upgraded 52 such profiles. We also extensively revised the metadata inconsistencies. Metadata fields, such as TF class and family, were revised and updated with the newest available information.

Together with motif collections, we provide additional analysis results, including motif clustering, familial TF binding motifs, predictions of TFBSs, and genome browser tracks. We clustered CORE collection motifs individually for each taxon and CORE and UNVALIDATED motifs. Clusters are visualized as radial trees and annotated with TF class information. Furthermore, we used the motif clusters to construct familial binding profiles. A familial binding motif represents DNA binding preference for closely related TFs. We used the latest versions of motifs in the CORE collection to scan the most popular genomes to predict TFBSs. Predicted binding sites are available to users along with the genome browser tracks. A new addition in 2022 was word clouds, generated for each TF to illustrate their functional roles. JASPAR database is user-friendly and accessible via the command line, RESTful API, R Bioconductor package, and pyJASPAR. We provide several computational tools to accompany the data. In this release, we added an enrichment analysis tool that allows the users to study the enrichment of the binding of a specific TF in a set of genomic regions and a TFBS extraction tool for easier extraction of TFBSs in a genomic region of interest.

My contributions to this paper included data collection and processing, manual curation of the motifs and their metadata, and revision of the existing motif metadata.

\section{Paper II. The release of the 2024 JASPAR database}

\textbf{\underline{Rauluseviciute I}}\footnotemark[\value{footnote}], Riudavets-Puig R\footnotemark[\value{footnote}], Blanc-Mathieu R, Castro-Mondragon JA, Ferenc K, Kumar V, Lemma RB, Lucas J, Chèneby J, Baranasic D, Khan A, Fornes O, Gundersen S, Johansen M, Hovig E, Lenhard B, Sandelin A, Wasserman WW, Parcy F, Mathelier A. \textit{JASPAR 2024: 20th anniversary of the open-access database of transcription factor binding profiles.} Nucleic Acids Research. 2024 Jan 5;52(D1):D174-82; doi: 10.1093/nar/gkad1059. \footnotetext{These authors should be regarded as joint first authors.}
\\

\textbf{Paper II} describes the latest release of the JASPAR database. In 2024, CORE and UNVALIDATED collections were enhanced or updated by 404 and 298 profiles, respectively. We validated and moved 72 profiles to the CORE collection. We expanded insects and nematodes motif collections by 90\% and 140\%, respectively. Plant TF classification was extensively revised and now is based on Plant-TFClass \cite{blanc2024plant}. One of the most significant updates in the 2024 release was the application of motif trimming. 75\% (1,869 out of 2,506) JASPAR 2022 motifs, and we trimmed all incoming new motifs. We applied an algorithm to remove the low IC positions on the sides of motifs to improve the quality of our collections. The trimming procedure will become a standard when preparing the motifs for manual curation in future releases.

The JASPAR database update process is complex, and we have refined it into three operations, abbreviated as "3D": \textbf{D}iscovery of the motifs, \textbf{D}ownstream analysis, and \textbf{D}eployment of the website. Each process has its software component as a workflow or a pipeline. The motif discovery operation finds new motifs in the peak data and processes them for manual curation. It also processes external motifs found in literature or databases. The downstream analysis prepares manually curated motifs for the website. In addition, the downstream analysis operation performs all additional analyses to produce the data needed for the website. Finally, the deployment workflow collects all produced data and updates the website. All software components received updates to implement new features, address users' needs, and increase the reliability and performance of the entire process of JASPAR update.

My contributions to this paper included data collection and processing, manual curation of the motifs and their metadata, and revision of the existing motifs. I contributed to the software development in the discovery and downstream analysis parts. I was responsible for motif trimming implementation and application and participated in the data analysis.

\section{Paper III. Identification of transcription factor co-binding patterns with non-negative matrix factorization}

\textbf{\underline{Rauluseviciute I}}, Launay T, Barzaghi G, Nikumbh S, Lenhard B, Krebs AR, Castro-Mondragon JA, Mathelier A. \textit{Identification of transcription factor co-binding patterns with non-negative matrix factorization.} bioRxiv. 2023:2023-04.
\\

\textbf{Paper III} describes a non-negative matrix factorization-based method to identify co-binding DNA patterns. The COBIND tool identifies DNA/RNA patterns in the neighborhoods of the sequences referred to as anchors. We applied COBIND to a collection of known transcription factor binding sites to automatically identify potential TF co-binding patterns in flanking regions. Interpreting such predicted DNA patterns can give insights into TF co-binding and the more complex gene regulation mechanism it creates.

In COBIND, NMF is applied to one-hot encoded sequences and pinpoints the enriched DNA patterns at fixed instances. We used 5,699 TFBS datasets from the UniBind database for 401 TFs in seven species \cite{puig2021unibind}. As a proof-of-concept, we present a classical example of POU5F1 co-binding with SOX2 or SOX17. For SOX2 and SOX17 TFBS datasets, we were able to recover the DNA pattern that represents the POU5F1 motif with different spacings, corresponding to canonical (SOX2::POU5F1 pair) and compressed (SOX17::POU5F1 pair) motifs. Through motif similarity and the protein-protein interaction (PPI) data, we inferred that POU5F1 most likely binds the discovered motif.

We evaluated all discovered co-binding patterns with motif similarity and PPI analysis. 78\% (462 out of 591) of the patterns had support in the PPI data. For example, TEAD1 is co-binding with TEAD family members. However, some inferred pairs have not yet been reported in the PPI data or do not have a good motif match. In \textit{A. thaliana} genome, the same co-binding pattern was found for HY5, ABF1, ABF3, ABF4, GBF2, and GBF3 anchors. Motif analysis revealed the similarity with the motif recognized by human NF-Y TFs (NF-YC and NF-YA). PPI data does not currently support these interactions. Another example where the PPI data did not support COBIND prediction is a co-binding pattern for CTCF. We found a co-binding pattern with two different spacings in \textit{H. sapiens}, \textit{M. musculus} and \textit{D. rerio} TFBS datasets. The literature reported that another TF does not bind the predicted DNA pattern but by CTCF itself. CTCF is a zinc-finger protein, and in a fraction of its binding sites, 9-11 zinc fingers bind to an extended portion of the motif.

We also demonstrated that the regions with predicted co-binding patterns are likely functionally relevant as they are more evolutionarily conserved than those with only the anchor TFBSs. Across species and all predictions, we systematically compared the conservation at the anchor locations and predicted co-binding sites. We observed that 23--83\% of the datasets from different species with predicted co-binding patterns showed significantly higher conservation at the co-binding motif locations compared to the alternative locations in regions only with the anchor motifs. For anchor sites, 31--67\% of the datasets demonstrated a more conserved anchor site in the genomic regions with predicted co-binding patterns. We observed for 23--64\% of datasets that the locations of both the anchor and co-binding patterns were more conserved in the co-bound regions.

Furthermore, we investigated regions with and without predicted co-binders in the open chromatin data by analyzing \ac{DHS} footprints from matching human cell types. By performing a similar analysis to the conservation, we found that in 85\% of co-binding pattern and DHS dataset pairs, DHS footprints were significantly deeper at the locations of the co-binding patterns in regions with predicted co-binding motifs compared to the genomic regions without (expectation by chance was 20\%). We observed significantly more profound DHS footprints in anchor locations with predicted co-binding patterns for 78\% of the pairs (expectation by chance was 21\%). We did not observe significance in a random expectation when considering deeper footprints for anchor and co-binding pattern sites together. In actual observation, it was significant for 73\% of the pattern-dataset pairs.

Finally, \ac{SMF} data from mouse embryonic stem cells (mESCs) confirmed that the predicted co-binding events associated with a set of TFs are likely occurring on the same DNA molecules. We analyzed predictions for 17 anchor TFs and found co-binding patterns in 1,694 regions. For each of these genomic regions and regions with only anchor motifs, we determined the proportions of molecules in five states: accessible, occupied by nucleosomes, occupied at the anchor motif site, only occupied at the co-binding motif site, or occupied at both sites. Overall, we found that a significantly more significant fraction of molecules was co-occupied at both sites in the regions with COBIND-predicted co-binding motifs compared to the genomic regions only with the anchor motifs. When looking at the individual anchors, significance was observed for 8 out of 17 analyzed anchor TFs.

We introduced a computational framework COBIND (\url{https://bitbucket.org/CBGR/cobind_tool/src/main/}) for DNA motif pattern discovery \textit{de novo} in the vicinity of predetermined regions or a set of sequences. During extensive analysis, we produced a collection of predicted co-binding patterns in 7 species, which we provide to the community at \url{https://cbgr.bitbucket.io/COBIND_results_page/}. DNA patterns discovered in the fixed locations from known TFBSs reveal the possible cooperative events and can be studied to investigate more complex TF-DNA interactions.

My contributions to this paper included developing the method, performing the analysis, and preparing the COBIND tool and resource with the results.